%% RESUMO NA LINGUA VERNACULA (elemento obrigatorio -- NBR 6028)
%% As palavras-chave sao definidas em config/dados.tex.

\begin{resumo}
Modelos de linguagem de grande escala vêm sendo usados na administração
pública para classificar documentos, mas não eliminam o erro, e a sua adoção
exige distinguir os casos em que a classificação automática é confiável
daqueles em que a revisão humana continua necessária. Esta dissertação
investiga em que medida sinais de incerteza produzidos pelos próprios modelos
permitem decidir quando aceitar automaticamente uma classificação e quando
encaminhá-la a uma pessoa. A fundamentação apoia-se na teoria da
classificação seletiva, que troca cobertura por acurácia, e na literatura
sobre autoavaliação e calibração de modelos de linguagem. A pesquisa,
quantitativa e aplicada, compara a classificação proposta pelos modelos com
uma classificação oficial atribuída por especialistas a uma base do setor
público e avalia cada sinal por meio de curvas risco-cobertura. Espera-se
oferecer um critério verificável de encaminhamento à revisão humana,
coerente com a exigência de supervisão proporcional ao risco.
\end{resumo}
